% UBT-AI-PROVENANCE-BEGIN
% schema: ubt-ai-provenance/v1
% tier: B_machine_verified
% ai_assistance: disclosed
% human_review: machine-verification
% editorial_responsibility: Ing. David Jaroš
% policy: ../../AI_PROVENANCE.md
% notice: Machine-verified against named sources or verifiers; individual attestation is not claimed.
% UBT-AI-PROVENANCE-END

\ifdefined\INCLUDEMODE\else
\documentclass[10pt]{article}
\usepackage[a4paper,margin=2cm]{geometry}
\usepackage{fontspec}
\usepackage{amsmath,amssymb}
\usepackage[hidelinks]{hyperref}
\usepackage{xurl}
\emergencystretch=3em
\usepackage{ubtprovenance}
\UBTTier{B}
\renewcommand{\UBTProvenanceNotice}{\par\smallskip{\footnotesize Provenience AI --- úroveň B. Text vznikl s pomocí AI; strojově ověřen uvedeným ověřovačem; individuální lidské potvrzení netvrdíme.}\par\smallskip}
\title{Krok 1: opravený test U(1) pozadí}
\author{UBT}
\date{2026-10-31}
\begin{document}\maketitle\UBTProvenanceNotice
\fi


% BILINGUAL-UNIT: u1-neutrality.scope

\section{Rozsah a opravený status}

GAP-U1-NEUTRALITY: FORMER PROOF REFUTED; PHYSICAL BRIDGE OPEN. Dřívější argument v \path{canonical/qed_phi_const/step1_u1_protection.tex} nahradil deklarované levé působení fáze konjugací. Opravujeme tento argument a pokračujeme pouze na původním nosiči jediného bikvaternionu. Mikroskopická akce, fyzikální reprezentace náboje a omezené kladné spektrum fluktuací zůstávají neurčené. Nezavádíme nové pole, metriku ani identifikaci vesmíru.

% BILINGUAL-UNIT: u1-neutrality.action

\section{Náboj určuje působení, nikoli komutace}

Pro elektromagnetický kalibrační parametr nezávislý na psi má každý Fourierův koeficient tentýž deklarovaný levý náboj. Reálnost ani nulový Fourierův index nezajišťují invarianci nenulového pole. Centrální konjugace by ponechala invariantní CELOU algebru a nemohla by nabít pouze ostatní složky. Původní věta tedy nedokazuje neutralitu ani bezhmotný foton.

\[
\begin{gathered}\Theta'=e^{i\alpha}\Theta,\qquad \left.\frac{d}{d\alpha}(e^{i\alpha}v)\right|_{\alpha=0}=iv\ne0\quad(v\ne0),\\ e^{i\alpha}ve^{-i\alpha}=v.
\qquad \theta'_n=e^{i\alpha}\theta_n.\end{gathered}
\]

% BILINGUAL-UNIT: u1-neutrality.mass

\section{Podmíněný test hmotnosti kalibračního bosonu}

Předpokládejme kladný kanonicky normovaný kinetický člen komplexního skaláru, hermitovský náboj Q a konstantní pozadí v v přirozených jednotkách se signaturou (+---). Srovnání jeho kvadratického kalibračního členu s Procovou normalizací dává následující výsledek. Zde je v samotná hodnota komplexního pole, nikoli konvence reálné radiální amplitudy v dělené odmocninou ze dvou. Nenulové pozadí v deklarované reprezentaci s jednotkovým nábojem tedy nesplňuje test neutrálního pozadí. Jde o diagnostiku původního skalárního srovnávacího modelu, nikoli o odvození dosud neurčené fyzikální Hessovy matice UBT. Samotná kalibrační invariance akce nezaručuje bezhmotný kalibrační boson.

\[
\begin{gathered}
\mathcal L_{\rm kin}=(D_\mu\phi)^\dagger D^\mu\phi,\qquad D_\mu=\partial_\mu+ieQA_\mu,\\
\mathcal L_{A^2}=e^2 v^\dagger Q^2v A_\mu A^\mu=\tfrac12m_\gamma^2 A_\mu A^\mu,\\
m_\gamma^2=2e^2\lVert Qv\rVert^2,\qquad m_\gamma^2=0\Longleftrightarrow Qv=0\quad(e\ne0).
\end{gathered}
\]

% BILINGUAL-UNIT: u1-neutrality.verification

\section{Ověření a pokračování}

Analytická oprava a nezávislé kontroly jsou zaznamenány v párové výzkumné poznámce \path{u1_neutrality_composite_channel} a \path{reports/u1_neutrality_2026_10_31.json}. Ověřovač \path{tools/verify_u1_neutrality.py} kontroluje fázový protipříklad a normalizaci hmotnosti kladného skaláru. LEAN-PENDING: není nainstalován lean ani lake. Tím není odvozena fyzikální Hessova matice UBT ani kanál multiversa. Původní závěr o QED je OPEN, nikoli PROVED. Výchozím jazykem překladu je angličtina. Netvrdíme nezávislé významové schválení překladu ani autorské potvrzení.

\ifdefined\INCLUDEMODE\def\UBTFinish{}\else\def\UBTFinish{\end{document}}\fi
\UBTFinish
